% BHU PYQ -- Professional Exam Template
\documentclass[11pt,a4paper]{article}

\usepackage[a4paper,top=2.5cm,bottom=2.5cm,left=2.8cm,right=2.8cm,headheight=14pt]{geometry}
\usepackage{amsmath,amssymb}
\usepackage[T1]{fontenc}
\usepackage[utf8]{inputenc}
\usepackage{lmodern}
\usepackage{microtype}
\usepackage{fancyhdr}
\usepackage{enumitem}
\usepackage{listings}
\usepackage{hyperref}
\usepackage{lastpage}
\usepackage{parskip}

\hypersetup{colorlinks=true,urlcolor=black,linkcolor=black,
  pdftitle={GEN-301: Digital Logic and Processors -- BHU PYQ}}

\pagestyle{fancy}
\fancyhf{}
\renewcommand{\headrulewidth}{0.4pt}
\renewcommand{\footrulewidth}{0.4pt}
\fancyhead[L]{\small   \textit{Banaras Hindu University}}
\fancyhead[R]{\small   \textit{GEN-301: Digital Logic and Processors}}
\fancyfoot[C]{\small Page  \thepage\ of \pageref{LastPage}}

\lstset{
  basicstyle=\small tfamily,
  frame=single,
  breaklines=true,
  columns=flexible,
  keepspaces=true
}


\newlist{parts}{enumerate}{2}
\setlist[parts,1]{label=(\alph*),leftmargin=*,itemsep=4pt,topsep=3pt}
\setlist[parts,2]{label=(\roman*),leftmargin=*,itemsep=2pt,topsep=2pt}

\newcommand{\pts}[1]{\hfill   \textit{[#1]}}


\begin{document}


\begin{center}
  {\large\bfseries BANARAS HINDU UNIVERSITY}\\[2pt]
  {\normalsize B.Voc. Semester III Examination, 2022-23}\\[6pt]
  \rule{0.8\linewidth}{0.5pt}\\[6pt]
  {\large\bfseries Paper: GEN-301 --- Digital Logic and Processors}\\[4pt]
  {\normalsize Time: Three hours \hfill Maximum Marks: 70}
\end{center}

\rule{\linewidth}{0.4pt}

\smallskip



\noindent   \textit{Instructions: Answer all Sections. Candidates are required to write their answers in their own words as far as practicable.}

\bigskip

\bigskip


\noindent {\large\bfseries SECTION --- A: Long Answer Questions}\pts{$15  \times 2 = 30$}
\rule{\linewidth}{0.4pt}\smallskip




\noindent   \textit{Instructions: Long answer type questions to be answered in about 750 words each. Attempt any two.}

\medskip


\noindent   \textbf{Question 1.} Prove the following Boolean identities:

\begin{parts}
  \item $A C + A B C = A C$
  \item $A + ar{A} B = A + B$
  \item $A B ar{C} + A ar{B} C + A B C = A(B + C)$
\end{parts}
\centerline{   \textbf{OR}}

\medskip


\noindent   \textbf{Question 2.} Explain about signal processing and discuss its different types in detail.

\medskip


\noindent   \textbf{Question 3.} Convert the number system from one form to another:

\begin{parts}
  \item $(6)_{10} = ( \quad )_2$
  \item $(110010)_2 = ( \quad )_8$
  \item $( \text{A9})_{16} = ( \quad )_{10}$
\end{parts}
\centerline{   \textbf{OR}}

\medskip


\noindent   \textbf{Question 4.} What are the basic gates in a logic circuit? Give a brief explanation of each with its respective truth table.

\bigskip


\noindent {\large\bfseries SECTION --- B: Short Answer Questions}\pts{$5  \times 6 = 30$}
\rule{\linewidth}{0.4pt}\smallskip




\noindent   \textit{Instructions: Attempt any six questions. Short answer type questions to be answered in about 550 words if theoretical.}

\medskip


\noindent   \textbf{Question 1.} What is a combinational circuit? Design a combinational circuit with two AND gates and one OR gate, and draw its truth table.

\medskip


\noindent   \textbf{Question 2.} Perform the following binary mathematical operations:

\begin{parts}
  \item Binary Addition:
    \[
    
\begin{array}{r@{\quad}l}
      1101_2 \\
    + \, 1010_2 \\
    \hline
    \end{array}
    \]
  \item Binary Subtraction:
    \[
    
\begin{array}{r@{\quad}l}
      1111_2 \\
    - \, 101_2 \\
    \hline
    \end{array}
    \]
\end{parts}

\medskip


\noindent   \textbf{Question 3.} What will be the output for the given logic circuit if $X_1 = 1$ and $X_2 = 0$?

\medskip


\noindent   \textbf{Question 4.} Calculate the number of cells in a Karnaugh Map (K-map) if we have 3 variables. Draw the K-map structure and write the cell representation.

\medskip


\noindent   \textbf{Question 5.} Simplify/Demorganize the following expression:
\[
\overline{(A+B)(C+D)}
\]

\medskip


\noindent   \textbf{Question 6.} Convert the given logical expression into Sum of Products (SOP) form:
\[
A B C + A ar{C} D + A B C D
\]

\medskip


\noindent   \textbf{Question 7.} What are latches and what are their types? Explain the NAND latch in detail.

\medskip


\noindent   \textbf{Question 8.} What is a multiplexer? Explain its working principles.

\bigskip


\noindent {\large\bfseries SECTION --- C: Very Short Answer Questions}\pts{$10  \times 1 = 10$}
\rule{\linewidth}{0.4pt}\smallskip




\noindent   \textit{Instructions: Answer the following questions in a single word or sentence.}


\begin{enumerate}[label=\arabic*.,leftmargin=*,itemsep=6pt]
  \item In K-map, what does the letter 'K' stand for? (Karnaugh)
  \item Give one example of a Universal Gate.
  \item What type of arithmetic operation is performed by an OR gate?
  \item "A binary variable can store the value of zero or one." Is this statement true or false?
  \item Which device converts a single input into multiple outputs? (Demultiplexer)
  \item Does a flip-flop store 1-bit of information?
  \item Is a register used for storing data?
  \item Multiply $(1001)_2$ with $(10)_2$.
  \item Write the full form of MSB. (Most Significant Bit)
  \item Write the 1's complement of $(1000)_2$.
\end{enumerate}

\vfill
\rule{\linewidth}{0.4pt}

\begin{center}\small
\href{https://bhu-exam.vercel.app/}{Click here to visit our website}\\[4pt]\href{https://me-aryan.vercel.app/}{Click here to visit our contributor}\\[4pt]\href{https://quashan-me.vercel.app/}{Click here to visit our contributor}
\end{center}

\end{document}
